\section{Implementation}
\label{sec:impl}

We implement \sysname as a CUDA C++ extension for PyTorch. The current backend
executes FP32 tensors in channels-last layout and provides the mask-aware
operators required by the evaluated networks: convolution, transposed
convolution, pooling, pointwise activation, addition, concatenation, and
upsampling. A model-conversion pass replaces compatible PyTorch modules while
preserving the graph topology, convolution parameters, and learned weights.
The upstream application continues to generate masks.

The construction kernel assigns contiguous output positions to each warp. A
thread tests one receptive field and terminates once it finds an active input.
A warp ballot identifies active lanes, population counts give their local
ranks, and one atomic operation reserves a contiguous worklist segment for the
warp. Active lanes then write their coordinates in lane order. This realizes
the tile-local ordering in Section~\ref{sec:design:construct} without a
device-wide prefix scan or sorting pass. An identity $1\times1$ convolution
has the same input and output coordinate space; when its input already carries
compatible mask and worklist metadata, the operator reuses that worklist and
skips construction.

The compute kernel uses a persistent, tiled implicit-GEMM organization. It
keeps partial sums in registers and stages input and weight tiles through
shared memory. Supported architectures use asynchronous global-to-shared
copies, with a synchronous path for earlier GPUs. Split-$K$ exposes additional
parallel work when the active-coordinate dimension alone cannot occupy the
device. Transposed convolution follows the same output-driven organization,
but partitions output coordinates by stride phase so that each phase visits
only its valid kernel offsets.

Efficient scheduling depends jointly on the GPU, layer shape, and activity.
We therefore autotune the sparse tile shape and split-$K$ factor and cache the
choice by device, convolution parameters, activity level, and execution
policy. For layers whose state-update semantics permit dense evaluation, the
candidate set can also include a cuDNN path. A sparse-only policy isolates the
worklist-driven kernel, and the cache key distinguishes the two policies and
the requested determinism mode. Autotuning and one-time weight preparation
complete before the timed sequence described in Section~\ref{sec:eval:setup}.

For fixed input values, weights, and mask, \sysname evaluates
Eq.~\eqref{eq:conv} at every active output coordinate. Floating-point
reordering may produce small numerical differences from cuDNN, so we verify
active outputs with numerical tolerances rather than requiring bitwise
identity. This check separates backend correctness from the task-quality
effect of the fixed upstream mask policy.
